\documentclass[10pt,a4paper]{amsart}
\usepackage[margin=19mm]{geometry}
\usepackage{amsmath,amssymb,mathtools}
\usepackage[scr=rsfs,cal=euler]{mathalfa}
\usepackage{xfrac}
\usepackage[unicode,hidelinks,bookmarks=false]{hyperref}
\newcommand{\ie}{i.e.\ }
\newcommand{\cf}{cf.\ }
\newcommand{\resp}{respectively\ }
\newcommand{\Subsection}{\subsection}
\providecommand{\nobreakdash}{\nobreak\hbox{-}}
\setlength{\emergencystretch}{2em}
\multlinegap0pt
\usepackage{amssymb}
\usepackage{tikz}
\usepackage[shortlabels]{enumitem}
\usepackage{bm}
\newtheorem{set}{Setting}
\newtheorem{rem}{Remark}
\newtheorem{prop}{Proposition}
\usepackage{cleveref}
\crefname{set}{Setting}{Settings}
\crefname{rem}{Remark}{Remarks}
\crefname{prop}{Proposition}{Propositions}
\def\K{\mathbb{K}}
\def\P{\mathbb{P}}
\def\bideg{\mathop{\rm bideg}}
\def\dim{\mathop{\rm dim}}
\title{Tensor product surfaces and graded syzygies}
\author{Matthew Weaver}
\date{Source: arXiv:2605.07974v1 (CC BY 4.0)}
\begin{document}
\maketitle

\noindent\emph{Source and licence.} Tensor product surfaces and graded syzygies. Version arXiv:2605.07974v1 (CC BY 4.0), distributed by arXiv under the Creative Commons Attribution 4.0 International licence, http://creativecommons.org/licenses/by/4.0/, as recorded in the arXiv licence metadata for this exact version and shown on the abstract page https://arxiv.org/abs/2605.07974v1. Full licence notice: \texttt{licenses/arxiv-2605.07974-CC-BY-4.0.txt}.\par\smallskip

\noindent\textit{Editorial scope.} Counted unit: the article's prop V (source0.tex lines 392--394). The article's complete original proof of it is delivered with it (source0.tex lines 397--401, one \texttt{proof} environment of 5 lines). No mathematics is written, repaired or re-derived: the statement and the proof are the article's own lines, in the article's own notation, and no retained line is substituted or re-flowed.  Retained uncounted as the source-local context the proof needs: lines 355--357; lines 364--366; lines 373--377.  Nothing was omitted from any retained span, so the package carries no omission record.  No retained line contains a citation, so the wrapper carries no bibliography and no bibliography entry.  No retained line contains a figure environment, an image-inclusion command or drawing code, so no image is delivered and the package declares no asset dependency.  Editorial formatting only, in addition: an AMS \texttt{amsart} preamble that restates the article's own declarations and macros used by the retained lines, namely the environments \texttt{set}, \texttt{rem}, \texttt{prop} on separate counters, together with the standard packages those lines need.  The retained environments print in this document as Set 1, Remark 1, Proposition 1 in order of appearance, whereas the article prints the same items with the numbering of its own full document.  The complete native source of the article as distributed in the arXiv e-print for this exact version is bundled unchanged as \texttt{sources/200124/source0.tex}.
\begin{set}\label{General Setting}
Let $R=\K[s,t,u,v]$ with $\bideg s,t =(1,0)$ and $\bideg u,v = (0,1)$. Let $U \subseteq R_{a,b}$ be a subspace with basis $\{p_0,p_1,p_2,p_3\}$ and let $I_U = (p_0,p_1,p_2,p_3) \subseteq R$. Assume that $U$ is basepoint free and let $\phi_U:\, \P^1\times \P^1 \longrightarrow \P^3$ be the regular map defined by $U$, with image $X_U$. Assume that $I_U$ has a first syzygy $S$ of bidegree $(0,n)$ and that $b\geq 2n-1$. Moreover, assume that $S$ is a singly graded syzygy in minimal degree, i.e. that $I_U$ has no syzygy in bidegree $(0,m)$ for $m<n$.
\end{set}

\begin{equation}\label{Syzygy rearrangement}
    0=\sum_{i=0}^3 \big(\sum_{j=0}^n a_{ij} u^{n-j}v^j\big) p_i = \sum_{j=0}^n \big( \sum_{i=0}^3 a_{ij} p_i\big) u^{n-j}v^j= \sum_{j=0}^n f_j u^{n-j}v^j
\end{equation}

\begin{rem}\label{f0 and fn nonzero}
We note that the set $\{f_0,\ldots,f_n\}$ is non-vanishing. Indeed, since $\{p_0,p_1,p_2,p_3\}$ is a basis of $U$, we note that $f_j =0$ if and only if $a_{0j} = a_{1j} = a_{2j} = a_{3j}=0$. As these are the coefficients of the polynomial entries of $S$, it follows that at least one of $f_0,\ldots,f_n$ is nonzero. In particular, from the assumptions above, it follows that both $f_0 \neq 0$ and $f_n \neq 0$. Indeed, if $f_0 =0$, then (\ref{Syzygy rearrangement}) shows that 
$$0=\sum_{j=1}^n f_j u^{n-j}v^j = v\cdot\sum_{j=1}^n f_j u^{n-j}v^{j-1}.$$
However, since $R$ is a domain, we then have $\sum_{j=1}^n f_j u^{n-j}v^{j-1} =0$. Hence by rearranging as in (\ref{Syzygy rearrangement}), it follows that $I_U$ has a syzygy in bidegree $(0,n-1)$, which is a contradiction. Similarly, it follows that $f_n\neq 0$.
\end{rem}

\begin{prop}\label{Dimension of V}
With the assumptions of \Cref{General Setting} and $V$ the subspace spanned by $\{f_0,\ldots,f_n\}$, we have that $2\leq \dim V\leq 4$.
\end{prop}

\begin{proof}
The latter inequality is clear, noting that $V$ is a subspace of $U$. Following \Cref{f0 and fn nonzero}, we have that $V \neq 0$, and so it suffices to show that $\dim V \neq 1$. Suppose that $\dim V =1$ and recall that $f_0 \neq 0$ by \Cref{f0 and fn nonzero}. Thus $\{f_0\}$ is a basis for $V$, and so we have $f_j = b_jf_0$ for some $b_j \in \K$, for $1\leq j\leq n$. Thus by (\ref{Syzygy rearrangement}), we see
$$0=\sum_{j=0}^n f_j u^{n-j}v^j = f_0 \cdot \sum_{j=0}^n b_j u^{n-j}v^j$$
where we take $b_0 =1$. Since $R$ is a domain and $f_0 \neq 0$, we have $\sum_{j=0}^n b_j u^{n-j}v^j=0$. However, this is impossible as $\{u^n,u^{n-1}v,\ldots,v^n\}$ is $\K$-linearly independent and $b_0\neq 0$. 
\end{proof}

\end{document}
